\section*{摘要：这篇论文研究什么}
接收器可以告诉我们哪些电流容易被观测，却不能直接告诉我们：由这些电流构成的模型会产生怎样的材料更新。本文从Maxwell方程出发，把一次指定电流改动拆成材料注入（material injection）、保留空间多重散射反馈（retained multiple-scattering feedback）和接收返回（receiver return），再追踪它怎样改变Gauss--Newton（GN）材料步。两组材料方向包含完整步差，但步差是否有用，还取决于曲率变化和当前残差的共同作用。

因此，我从接收基底出发，在固定电流维数下测试少量一列移出、一列移入的条件交换（conditional exchange）。更丰富的约化模型先排序，每轮最多两个候选接受完整定向切向作用$Jd$的数值复核；只有共同完整二次目标确实下降，才改变空间。接受阈值只使用在线数值尺度，不读取完整最优材料步或真值。原四对象的对象汇总完整GN gap中位下降65.2\%；冻结规则后，十个完整额外对象全部改善，中位下降50.7\%。另一个对象保留晚期前置求解失败。这里的gap衡量材料步对完整GN步的保真度，不是最终图像误差。暗电流、信息正而收益负、以及单方向负而联合方向正的物理见证，共同说明接收强度和信息汇总不足以决定当前材料用途。论文给出的结果是一个有明确费用和适用范围的电磁机制与条件设计。
